% =====================================================================
%  Capítulo 4. Análisis y diseño
% =====================================================================
\chapter[ANÁLISIS Y DISEÑO]{Análisis y diseño}
\label{cap:analisis-diseno}

El análisis consiste en identificar qué debe ocurrir dentro del sistema y cómo se relacionarán
entre sí los distintos procesos que lo conforman, lo que permite entender el problema y documentar
las relaciones entre los usuarios, el programa y los componentes que integran la lógica del
sistema antes de escribir cualquier línea de código \cite{pressman2010}. En el caso del presente
trabajo, esta etapa corresponde a las fases de definición de requerimientos y de diseño de la
arquitectura del modelo incremental descrito en el Capítulo~\ref{cap:propuesta}, y se orienta a
describir, mediante requerimientos y diagramas, el comportamiento esperado de la interfaz del P300
Speller, del módulo de detección de la componente P300 y del módulo de predicción de palabras,
combinando el lenguaje del usuario con el enfoque técnico de quien desarrolla el sistema.

Por otra parte, el diseño se encarga de traducir dicho análisis en decisiones concretas de
implementación: a partir de los requerimientos y los diagramas generados se define cómo se
organizará el código en módulos, cómo se intercambiará la información entre la interfaz, el
modelo de aprendizaje profundo y el modelo de lenguaje, qué datos deberán almacenarse, así como la
plataforma en la que se ejecutará cada componente y los recursos que requerirá \cite{pressman2010}.

\section[REQUERIMIENTOS]{Requerimientos}

%% ---------------------------------------------------------------------
%% NOTA: los siguientes requerimientos son una PROPUESTA inicial basada
%% en el protocolo. Revisar, eliminar o agregar según lo que se decida.
%% ---------------------------------------------------------------------

\subsection{Requerimientos funcionales}

\begin{itemize}
  \item \textbf{RF1:} El sistema deberá proveer una interfaz gráfica que muestre una matriz de
  selección de caracteres de $6\times6$, con letras, números y comandos básicos de edición.
  \item \textbf{RF2:} El sistema deberá intensificar de manera aleatoria las filas y columnas de la
  matriz durante cada secuencia de estimulación.
  \item \textbf{RF3:} El sistema permitirá configurar los parámetros de estimulación, como la
  duración de la intensificación, el intervalo entre estímulos y el número de repeticiones por
  selección.
  \item \textbf{RF4:} El sistema deberá cargar registros de EEG provenientes de bases de datos
  públicas del paradigma P300 Speller, junto con la información de los estímulos asociados.
  \item \textbf{RF5:} El sistema deberá preprocesar las señales EEG mediante filtrado pasa banda,
  remuestreo, segmentación en épocas y normalización por canal.
  \item \textbf{RF6:} El sistema utilizará un modelo de aprendizaje profundo para clasificar cada
  época como objetivo o no objetivo, asignándole un puntaje de probabilidad.
  \item \textbf{RF7:} El sistema deberá determinar el carácter seleccionado a partir de la
  acumulación de los puntajes de cada fila y cada columna a lo largo de las repeticiones.
  \item \textbf{RF8:} El sistema mostrará en la interfaz el texto construido por el usuario,
  actualizándolo después de cada selección.
  \item \textbf{RF9:} El sistema contará con un módulo de predicción de palabras que genere
  sugerencias a partir de los caracteres escritos hasta ese momento.
  \item \textbf{RF10:} El sistema permitirá seleccionar una palabra sugerida mediante el mismo
  mecanismo de estimulación y detección de la componente P300.
  \item \textbf{RF11:} El sistema permitirá borrar el último carácter o palabra seleccionados
  mediante un comando de la matriz.
  \item \textbf{RF12:} El sistema permitirá configurar el subconjunto de canales EEG empleados por
  el modelo de detección.
  \item \textbf{RF13:} El sistema mostrará al finalizar una sesión las métricas de desempeño
  obtenidas, como la exactitud por carácter, el tiempo por selección y la tasa de transferencia de
  información.
\end{itemize}

\subsection{Requerimientos no funcionales}

\begin{itemize}
  \item \textbf{RNF1:} El sistema deberá implementarse en el lenguaje de programación Python.
  \item \textbf{RNF2:} El sistema deberá ejecutarse en una computadora con sistema operativo
  Windows.
  \item \textbf{RNF3:} La presentación de los estímulos deberá sincronizarse con los instantes
  registrados de cada intensificación, de modo que cada época corresponda al estímulo que la
  provocó.
  \item \textbf{RNF4:} La clasificación de las épocas de una selección deberá completarse en un
  tiempo que no retrase de manera perceptible la presentación de la siguiente secuencia de
  estímulos.
  \item \textbf{RNF5:} La interfaz deberá presentar la matriz con un tamaño de carácter y un
  contraste suficientes para que el usuario pueda fijar su atención sin esfuerzo visual excesivo.
  \item \textbf{RNF6:} El sistema deberá tener una arquitectura modular que permita reemplazar el
  modelo de detección o el modelo de predicción de palabras sin modificar la interfaz.
  \item \textbf{RNF7:} La etapa de inferencia del modelo deberá poder evaluarse en un
  microordenador NVIDIA Jetson.
  \item \textbf{RNF8:} El código fuente deberá estar documentado y acompañarse de un manual de
  usuario que describa la instalación, la ejecución y el uso general del sistema.
\end{itemize}

\section[CASOS DE USO]{Casos de uso}

Los casos de uso representan las distintas funcionalidades del sistema y se definen a partir de los
requerimientos funcionales, de manera que cada uno describe una acción específica, junto con sus
precondiciones, pasos y posibles errores, que el sistema debe realizar para cumplir con las
necesidades de quienes interactúan con él \cite{pressman2010}. Los actores, por su parte, son los
roles que se asignan a las personas o a los sistemas externos que interactúan con dichas
funcionalidades; en el presente trabajo intervienen principalmente el usuario que construye
mensajes a través del P300 Speller y quien configura los parámetros de estimulación y del modelo
de detección.

%% TODO: tabla de actores, tablas de casos de uso y diagramas de casos de uso.
